\documentclass[../../../include/open-logic-section]{subfiles}

\begin{document}

\olfileid{sol}{met}{int}

\olsection{引言}

一阶逻辑具有若干良好性质。我们知道它是不可判定的，但它至少是可公理化的。
也就是说，存在一阶逻辑的证明系统，它们是可靠且完备的；
即，它们产生可推导关系~$\Proves$，
满足：对任意语句集~$\Gamma$ 和语句~$!A$，$\Gamma \Entails !A$ 当且仅当 $\Gamma \Proves !A$。
这尤其意味着一阶逻辑的有效式是可计算枚举的。% Part: second-order-logic
存在一个可计算函数~$f\colon \Nat \to \Sent[L]$，
使得~$f$ 的值恰好是且仅是~$\Lang{L}$ 的所有有效语句。
这是因为推导本身可以被枚举，而那些推导出单一语句的推导则可以映射到该语句。
二阶逻辑比一阶逻辑表达能力更强，因此一般而言，刻画它的有效式也更为复杂。
事实上，我们将会证明，二阶逻辑不仅是不可判定的，甚至它的有效式也不是可计算枚举的。
这意味着二阶逻辑不可能有可靠且完备的证明系统（尽管可靠但不完备的证明系统是存在的，并且事实上是重要的研究对象）。% Chapter: metatheory

一阶逻辑还有另外两个性质：它是紧致的（如果一个语句集~$\Gamma$ 的每个有穷子集都是可满足的，那么~$\Gamma$ 本身也是可满足的），
并且 勒文海姆--斯科伦 定理对它成立（如果 $\Gamma$ 有一个无穷模型，那么它也有一个可数无穷模型）。
这两个结果对于二阶逻辑都不成立。原因同样在于，二阶逻辑能够表达关于论域大小的事实，而一阶逻辑则不能。% Section: introduction

\end{document}